\chapter{Evaluation}

Evaluating an agentic system requires choosing, at the outset, what the evaluation is 
intended to measure. The simplest approach is outcome-based evaluation: given a task, did 
the agent produce the correct final answer? Its limitation is that a correct answer carries 
no information about the path taken to reach it. An agent may succeed on a given instance 
through a sequence of steps that is inefficient, brittle, or unlikely to generalise.

Trajectory-based evaluation addresses this by treating the sequence of decisions as the 
object of study, not merely its outcome. Rather than scoring the final answer alone, the 
evaluator inspects the full sequence of tool calls, retrieval steps, and intermediate 
reasoning: whether the right tools were selected, whether they were invoked with valid 
arguments, and whether each step was necessary given the state of the task at that point. 
This approach is more demanding to instrument and score, but it produces a more diagnostic 
picture of system behaviour, one that can distinguish a robust agent from one that happens 
to succeed on the test distribution.

\section{Traces}

The practical foundation of both evaluation modes is the trace: a structured, replayable 
record of a single agent execution. A trace captures the user input, every model call, every 
tool invocation and its result, any retrieved context, the final answer, per-step latency, 
and accumulated cost. Without traces, evaluation is limited to comparing inputs and outputs; 
with them, it becomes possible to isolate the specific decision that caused a failure.

Traces must reflect the hierarchical structure of agent execution. When a parent agent 
delegates to a sub-agent, that nested execution should appear as a child span within the 
parent trace. When a tool call fails, the error should be attached to the corresponding 
observation rather than surfaced only at the top level. This structure is what makes a trace 
analytically useful rather than merely a log.

\section{Scores}

Evaluation results are represented as scores attached to traces or to individual spans within 
them. Scores can be produced by automated checks, by human annotators, or by LLM-based 
judges, and they can target different levels of the execution. At the outcome level, the 
relevant dimensions include whether the agent satisfied the user's objective, whether factual 
claims are correct, and whether assertions are grounded in retrieved evidence or tool outputs. 
At the trajectory level, relevant dimensions include whether the agent selected appropriate 
tools, whether arguments were valid and complete, and whether tool results were correctly 
incorporated into subsequent reasoning. Orthogonal to both are efficiency scores, measuring whether 
the task was completed within acceptable cost and latency budgets, and safety scores, which 
assess whether the agent respected applicable permissions and constraints.

\section{Datasets and Experiments}

Langfuse supports input/expected-output datasets as the baseline format for offline 
evaluation. As the system is deployed and production traces accumulate, the dataset can be 
versioned and progressively enriched: traces flagged during online audit, where the agent 
failed or behaved unexpectedly, become new labelled examples in the next dataset version. 
This mechanism ensures that the evaluation suite tracks the actual failure modes encountered 
in production, and that fixes introduced in a new agent version are verified against the 
exact conditions that exposed the original defect.

\section{Stratified Sampling with Unbalanced Weights}

Production traffic is rarely uniform. In a typical deployment, the majority of conversations 
receive no explicit feedback, a small fraction receive a positive signal such as a thumbs-up, 
and a smaller fraction still receive a negative signal. Annotating a simple random sample of 
this traffic would leave the most informative cases, especially those where the agent visibly 
failed, severely underrepresented.

Stratified sampling addresses this by partitioning traffic into buckets and allocating 
annotation capacity to each bucket independently. In the feedback example, the natural 
partition is:

\begin{center}
\begin{tabular}{ll}
\toprule
\textbf{Bucket} & \textbf{Meaning} \\
\midrule
$h = \uparrow$   & Thumbs up \\
$h = \downarrow$ & Thumbs down \\
$h = \emptyset$  & No feedback \\
\bottomrule
\end{tabular}
\end{center}

Let $W_h$ denote the true production share of bucket $h$, and let $n_h$ denote the number 
of examples annotated within that bucket. For each bucket, the observed defect rate is:

\[
\hat{p}_h = \frac{\text{number of defective examples in bucket } h}{n_h}
\]

The overall production defect rate is then recovered by re-weighting each bucket estimate 
by its production share:

\[
\hat{p}_{\text{overall}} = \sum_h W_h \hat{p}_h
\]

This is the stratified sampling correction: because annotation overrepresents certain 
buckets, each bucket-level estimate must be scaled back to its true weight before 
aggregation. The approximate variance of this estimator is:

\[
\mathrm{Var}(\hat{p}_{\text{overall}}) = \sum_h W_h^2 \, \frac{\hat{p}_h(1 - \hat{p}_h)}{n_h}
\]

and the resulting 95\% confidence interval is:

\[
\hat{p}_{\text{overall}} \pm 1.96 \sqrt{\sum_h W_h^2 \, \frac{\hat{p}_h(1 - \hat{p}_h)}{n_h}}
\]

A practical consequence of this formula is that uncertainty in the overall estimate is 
dominated by the largest production bucket. If the no-feedback population represents 98.5\% 
of traffic, then $W_{\emptyset}^2$ is an order of magnitude larger than the weights of the 
other buckets, and the no-feedback variance term will dominate the confidence interval 
regardless of how many thumbs-down examples are annotated. Annotating high-signal buckets 
such as thumbs-down is valuable for diagnosing failure modes, but it does not substitute 
for adequate coverage of the no-feedback population when the goal is to estimate overall 
production performance reliably.

\subsection*{Numerical example}

Suppose production traffic is distributed as follows:

\begin{center}
\begin{tabular}{lrrr}
\toprule
\textbf{Bucket} & \textbf{Prod.\ share} $W_h$ & \textbf{Annotated} $n_h$ 
    & \textbf{Defect rate} $\hat{p}_h$ \\
\midrule
Thumbs up   & 1.0\%  & 200   & 2\%  \\
Thumbs down & 0.5\%  & 500   & 65\% \\
No feedback & 98.5\% & 2\,000 & 8\%  \\
\bottomrule
\end{tabular}
\end{center}

The weighted defect estimate is:

\[
\hat{p}_{\text{overall}} = 0.010 \times 0.02 + 0.005 \times 0.65 + 0.985 \times 0.08 
= 0.0823
\]

The standard error is:

\[
SE = \sqrt{
  0.010^2 \cdot \frac{0.02 \times 0.98}{200}
  + 0.005^2 \cdot \frac{0.65 \times 0.35}{500}
  + 0.985^2 \cdot \frac{0.08 \times 0.92}{2000}
} \approx 0.00598
\]

The 95\% confidence interval is therefore $[7.05\%,\ 9.40\%]$. The no-feedback term 
contributes virtually all of the variance, illustrating that increasing the thumbs-down 
annotation budget would narrow the interval only marginally.

\section{LLM-as-a-Judge}

Some evaluation criteria resist rule-based scoring: helpfulness, clarity, groundedness, and 
policy compliance each require a degree of interpretive judgment. LLM-based judges can 
operationalise these criteria at scale, but they introduce systematic biases that must be 
accounted for, including verbosity bias, position bias, and a tendency toward self-preference 
when the judge and the agent under evaluation share the same base model. Mitigation requires 
explicit, criterion-anchored rubrics; eliciting a chain of evidence before requesting a 
numerical score; enforcing structured output; running at low temperature; and calibrating 
judge scores against human annotations on a held-out subset.

LLM judges should not displace deterministic checks where deterministic checks are 
applicable. Conformance to a JSON schema should be verified by a schema validator. 
Functional correctness of generated code should be verified by executing tests. The role 
of an LLM judge is to handle the residual criteria that deterministic methods cannot reach.

\section{Failure Taxonomy}

The diagnostic value of traces is substantially higher when failures are categorised 
consistently. A working taxonomy for agentic systems covers three failure modes: incorrect 
tool usage, which encompasses both wrong tool selection and malformed input parameters; 
incorrect output, which covers both hallucination and wrong format; and guardrails 
violation. Labelling failures according to such a taxonomy transforms a collection of 
individual errors into a distribution of failure modes, which in turn makes it possible to 
prioritise improvements and measure whether interventions are effective.
\section{Offline and Online Evaluation}

Offline evaluation runs a fixed dataset before deployment and functions as a gate: a 
promotion is blocked if quality falls below a defined threshold. Online evaluation samples 
live production traffic after deployment and functions as a monitor: it detects distribution 
shift, novel user behaviour, tool failures, and model drift that fixed datasets do not 
anticipate. The two modes are complementary and a mature system requires both. Offline 
evaluation gives deployment confidence; online evaluation gives operational visibility. 
Human annotation periodically calibrates the automated scoring pipeline, and production 
trace mining continuously supplies new candidate items for the offline dataset, closing the 
loop between the two modes.

\section{Summary}

Agent evaluation is not a single procedure but a system of interlocking practices: choosing 
between outcome and trajectory evaluation, instrumenting traces, defining scored dimensions, 
maintaining a versioned dataset, deploying LLM judges with appropriate mitigations, and 
operating both offline and online pipelines. The connecting principle is that every 
production failure should become a labelled example, every labelled example should become 
an evaluation, and every evaluation should become a deployment gate.